\documentclass[11pt,a4paper,english,twoside]{article}

\usepackage{a4wide}
\usepackage{amssymb, amsmath, amsthm}
\usepackage{graphicx}
\usepackage{subcaption}
\usepackage[all]{xy}
\usepackage[pdftex,hyperref,svgnames]{xcolor}
\usepackage[pdftex,colorlinks=true,
pdfstartview=FitV,
pdfnewwindow=true,
linktoc = page,
linkcolor= blue,
citecolor= red,
urlcolor= blue,
hyperindex=true,
hyperfigures=false]{hyperref}
\hypersetup{linktocpage}
\usepackage{dsfont}
\usepackage{empheq}
\usepackage[nosort, nocompress]{cite}
\usepackage{float}
\usepackage{cancel}
\usepackage{relsize}
\usepackage{soul}
\usepackage{enumerate}
\usepackage{enumitem}
\usepackage{hhline}
\usepackage{multirow}
\usepackage{xspace}
\usepackage{bbm}
\usepackage{pdfpages}
\usepackage{setspace}
\usepackage{comment}

\usepackage{tocloft}
\renewcommand\cftsubsecafterpnum{\vskip4pt}
\renewcommand\cftsubsubsecafterpnum{\vskip1pt}
\setlength{\cftbeforesubsecskip}{6pt}
\setlength{\cftbeforesecskip}{11pt}


\newcommand{\nn}{\nonumber}

\def\beq{\begin{equation}}
\def\eeq{\end{equation}}
\def\bea{\begin{eqnarray}}
\def\eea{\end{eqnarray}}

\def\del {\partial}
\def\d {{\rm d}}
\def\w {\wedge}



%******************************************

\begin{document}
\numberwithin{equation}{section}


%\tableofcontents

%\newpage

\section{Model conventions}

We consider single-field 4d cosmological models of the form
\beq
{\cal S} = \int \d^4 x \sqrt{|g_4|} \left( \frac{M_p^2}{2} {\cal R}_4 - \frac12 (\del \varphi)^2 - V(\varphi) - {\cal L}_m - {\cal L}_r \right) \ ,
\eeq
following conventions of \cite[Sec. 3.2]{Andriot:2026lac}. We allow for coupling functions $A_m(\varphi),\, A_r(\varphi)$ to matter and radiation, in the corresponding Lagrangians, in such a way that
\beq
\rho_m = A_m(\varphi) \, \bar{\rho}_m \ ,\quad \rho_r = A_r(\varphi) \, \bar{\rho}_r \ , \quad {\rm where} \ \bar{\rho}_m = \bar{\rho}_{m0} \, a^{-3} \ ,\ \bar{\rho}_r = \bar{\rho}_{r0} \, a^{-4}\ . 
\eeq
We are going to look for cosmological solutions with FLRW metric. We will take the following fiducial values today
\bea
{\rm Today}:&& \quad \Omega_{r0} = 0.0001 \ ,\ \Omega_{m0}= 0.3149 \ ,\ \Omega_{\varphi0}=0.6850 \ ,\ \Omega_{k0}=0 \ ,\\
&& \quad A_r(\varphi_0)=1 \ ,\ A_m(\varphi_0)=1 \ ,\nn
\eea
and we will fix $w_{\varphi0}$ to get realistic solutions.

As argued in \cite{Andriot:2025los}, such a setting allows to define an effective dark energy, whose energy density and equation of state parameter are given by
\beq
\rho_{{\rm DE}} = \rho_{\varphi} + (A_m - 1) \bar{\rho}_m + (A_r - 1) \bar{\rho}_r \ ,\ w_{{\rm DE}} = \frac{w_{\varphi} + \frac13 (A_r -1) \frac{\bar{\rho}_r}{\rho_\varphi}}{1 + (A_m-1) \frac{\bar{\rho}_m}{\rho_\varphi} + (A_r-1) \frac{\bar{\rho}_r}{\rho_\varphi}  } \ .
\eeq
The coupling to matter can provide an effective phantom dark energy in the recent universe.

\section{Tackling the Hubble tension}

As pointed-out in \cite{Andriot:2025los}, the realistic solutions with a coupling to matter exhibit a systematic bump in $\Omega_{{\rm DE}}$ slightly after radiation - matter equality. This is reminiscent of Early Dark Energy (EDE) and might help for the Hubble tension. After a few tests, for example with the models described in \cite[Fig. 6]{Andriot:2026lac}, the resulting change in $H$ is however minor. The situation could be improved with an earlier bump. The latter could be generated by a coupling to radiation (also thanks to the earlier kination phase). We thus consider the following two models as examples.
\bea
&&\text{Model and solution 1:} \label{model1}\\
&& V(\varphi)= V_0 \, e^{- \sqrt{\frac83}\, \varphi} \ ,\ A_m(\varphi)= 1 - \frac{1}{16} \left(1-e^{\sqrt{\frac38}\, \varphi}\right) \ ,\ A_r(\varphi) = 1 + \frac{1}{16} \left(1- e^{\sqrt{\frac32}\, \varphi} \right)\nn\\
&& w_{\varphi0} = -0.59420916809 \ ,\ \varphi_0 =0 \ ,\nn
\eea
and
\bea
&&\text{Model and solution 2:} \label{model2}\\
&& V(\varphi)= V_0 \, e^{- \sqrt{\frac83}\, \varphi} \ ,\ A_m(\varphi)= 1 - \frac{1}{16} \left(1-e^{\sqrt{\frac38}\, \varphi}\right) \ ,\ A_r(\varphi) = 1 - \frac{1}{50} \left(1- e^{-1\, \varphi} \right)\nn\\
&& w_{\varphi0} = -0.59422361340 \ ,\ \varphi_0 =0 \ ,\nn
\eea
where $V_0$ and $\dot{\varphi}_0$ are fixed by $\Omega_{\varphi0},\, w_{\varphi0}$. We display the solutions in Figure \ref{fig:models}.
\begin{figure}[ht!]
\begin{center}
\begin{subfigure}[H]{0.48\textwidth}
\includegraphics[width=\textwidth]{figmodel1O.pdf}\caption{$\bar{\Omega}_r, \bar{\Omega}_m, \Omega_{{\rm DE}}$}\label{fig:model1O}
\end{subfigure}\quad
\begin{subfigure}[H]{0.48\textwidth}
\includegraphics[width=\textwidth]{figmodel2O.pdf}\caption{$\bar{\Omega}_r, \bar{\Omega}_m, \Omega_{{\rm DE}}$}\label{fig:model2O}
\end{subfigure}\\
\begin{subfigure}[H]{0.48\textwidth}
\includegraphics[width=\textwidth]{figmodel1logrho.pdf}\caption{}\label{fig:model1logrho}
\end{subfigure}\quad
\begin{subfigure}[H]{0.48\textwidth}
\includegraphics[width=\textwidth]{figmodel2logrho.pdf}\caption{}\label{fig:model2logrho}
\end{subfigure}\\
\begin{subfigure}[H]{0.48\textwidth}
\includegraphics[width=\textwidth]{figmodel1wDE.pdf}\caption{}\label{fig:model1wDE}
\end{subfigure}\quad
\begin{subfigure}[H]{0.48\textwidth}
\includegraphics[width=\textwidth]{figmodel2wDE.pdf}\caption{}\label{fig:model2wDE}
\end{subfigure}
\caption{Solution for model 1 \eqref{model1} (left column) and model 2 \eqref{model2} (right column). We also display the resulting $w_{{\rm DE}}$ in the recent universe (green curve), compared to DESI + CMB + Union 3 with CPL parametrisation (orange) from \cite{DESI:2025zgx}.}\label{fig:models}
\end{center}
\end{figure}

We definitely see an early ``bump'' in $\Omega_{{\rm DE}}$ in Figure \ref{fig:model1O} and \ref{fig:model2O}. It is unclear to us if this can really accommodate the Hubble tension, because there are a priori differences with the standard EDE model. In particular, there is no initial rise in the ``bump'', we are not sure the plateau shape and duration is fine, and we are not sure whether it decreases fast enough.

Note that $w_{{\rm DE}}$ in the recent universe exhibits phantom crossing as expected. The fit with DESI data could be improved by adjusting the matter coupling.


\section{Energy-momentum tensor and perturbations}

To implement perturbations in CLASS, it is useful to rederive the energy-momentum tensor and perturb it. To that end, we follow conventions from \cite[App. B]{Andriot:2026lac}, namely that
\beq
{\cal S} = \int \d^4 x \left(\frac{M_p^2}{2} {\cal R}_4 - {\cal L} \right) \ ,\quad T_{\mu\nu} = -g_{\mu\nu} {\cal L} + 2 \frac{\delta {\cal L}}{\delta g^{\mu\nu}} \ .
\eeq
We will use that both the background and the perturbed metric are block diagonal between time (index $0$) and space (index $i,j,k$). This is all we will need for now. We also recall that $T_{00} = \rho$, $T_{ij} =g_{ij}\, p$.

We perform the computation for
\bea
&& {\cal L} = \frac{1}{2} (\del \varphi)^2 + V(\varphi) + {\cal L}_m + {\cal L}_r \\
&& {\cal L}_{m/r}  = A_{m/r}(\varphi)\ \bar{{\cal L}}_{m/r} \ ,\quad \bar{{\cal L}}_{m} = \bar{{\cal L}}_{m0}/ \sqrt{|g_3|} \ ,\quad \bar{{\cal L}}_{r} = \frac12 F_{\mu\nu} F^{\mu\nu} \ ,
\eea
where $|g_3|$ denotes the determinant of the spatial metric block, and $(\del \varphi)^2 = \del^{\mu} \varphi \del_{\mu} \varphi$. We obtain
\beq
T^{\mu}{}_{\nu} = \del^{\mu} \varphi \del_{\mu} \varphi - \delta^{\mu}_{\nu} \frac12 (\del \varphi)^2 - \delta^{\mu}_{\nu} V + {\cal L}_m \left( - \delta^{\mu}_{\nu} + \delta^{\mu}_{i} \delta^{i}_{\nu} \right)  + {\cal L}_r  \left( - \delta^{\mu}_{\nu} + \frac43 \delta^{\mu}_{i} \delta^{i}_{\nu} \right) \ . 
\eeq
To get the radiation contribution, we made several assumptions, which might be justified as gauge choices and by isotropy and homogeneity. For example, it is convenient to assume that there is no electric field, $F_{0i}=0$, and that all magnetic field components are equal, in a space basis where all metric components are equal and diagonal. These choices however appear problematic with respect to the shear $\Sigma^i{}_j$, later. Note that the correct $\rho_r, p_r$ are still obtained this way. More explicitly, for any overall normalization $c$ in $\bar{{\cal L}}_r=c F_{\mu\nu}F^{\mu\nu}$, the isotropic magnetic configuration gives the radiation tensor $T^0{}_0=-{\cal L}_r$, $T^i{}_j=\frac13{\cal L}_r\delta^i_j$; the value of $c$ only fixes the normalization of the field amplitude inside $\bar\rho_r$. In the following we use the convention above and identify $\rho_m= {\cal L}_m$, $\rho_r= {\cal L}_r$, and the same for the bar quantities. For the shear, we get
\beq
\Sigma^i{}_j = T^i{}_j - \frac13 \delta^i_j\, T^k{}_k = \del^i \varphi \del_j \varphi - \frac13 \delta^i_j \,  \del^k \varphi \del_k \varphi  \ .
\eeq

We now turn to the perturbation. Beyond the block structure of the metric mentioned above, we add that background scalar field is only time dependent, but its perturbation $\delta \varphi$, can be space dependent as well. We then obtain at linear order, with $V'=\del_{\varphi} V$,
\begin{align}
\delta\Sigma^{i}{}_{j} = &\ 0 \ ,\\[4pt]
\delta T^{0}{}_{0}
= &\ \frac12\dot{\varphi}^2\,\delta g^{00}
   + g^{00}\dot{\varphi}\,\delta\dot{\varphi}
   - V'\delta\varphi
   - \delta\bar{{\cal L}}_m - \delta\bar{{\cal L}}_r \\
& \ - (A_m -1)\,\delta\bar{{\cal L}}_m - (A_r-1)\,\delta\bar{{\cal L}}_r
   - \delta\varphi\, \rho_{m}\, \del_{\varphi} \ln A_m 
   - \delta\varphi\, \rho_{r}\, \del_{\varphi} \ln A_r \ ,\nn
\\[4pt]
\delta T^{i}{}_{0} = &\ g^{ij}\,\dot{\varphi}\,\partial_j\delta\varphi \ ,
\\[4pt]
\delta T^{i}{}_{j}
= &\ -\delta^{i}_{j}\, \frac{1}{2} \dot\varphi^{2}\,\delta g^{00}
   - \delta^{i}_{j}\,g^{00}\dot{\varphi}\,\delta\dot{\varphi}
   - \delta^{i}_{j}\,V'\, \delta\varphi + \frac13\delta^{i}_{j}\, \delta\bar{{\cal L}}_r \\
&\ + \frac13\delta^{i}_{j}\left((A_r-1)\,\delta\bar{{\cal L}}_r
   +  \delta\varphi\, \rho_{r}\, \del_{\varphi} \ln A_r \right) \ ,
\end{align}   
where we put on second lines the new terms due to the couplings.

We further obtain
\begin{align}
\nabla_{\mu}\delta T^{\mu}{}_{0}
= &\ \nabla^i \left(\dot{\varphi}\,\del_i \delta\varphi \right) + \nabla_{0}\left( \frac12\dot{\varphi}^2\,\delta g^{00}
   + g^{00}\dot{\varphi}\,\delta\dot{\varphi}
   - V'\delta\varphi
   - \delta\bar{{\cal L}}_m - \delta\bar{{\cal L}}_r  \right)\\
&\ - (A_m -1)\, \nabla_{0}\delta\bar{{\cal L}}_m - (A_r-1)\, \nabla_{0}\delta\bar{{\cal L}}_r - \del_{\varphi} A_m \, \dot{\varphi}\, \delta\bar{{\cal L}}_m - \del_{\varphi} A_r \, \dot{\varphi}\, \delta\bar{{\cal L}}_r   \nn\\
&\ - \left(\rho_m\, \del_{\varphi} \ln A_m +\rho_r\, \del_{\varphi} \ln A_r\right)\delta\dot{\varphi}
   - \left(\rho_m \frac{\del_{\varphi}^2 A_m}{A_m} + \rho_r \frac{\del_{\varphi}^2 A_r}{A_r} \right)\dot{\varphi}\,\delta\varphi \ ,\nn
\\[4pt]
\nabla_{\mu}\delta T^{\mu}{}_{j}
= &\ g^{00}\ddot\varphi\,\partial_j\delta\varphi
   + g^{00}\dot\varphi\,\nabla_0\partial_j\delta\varphi
   - \nabla_j\!\left(\frac12\delta g^{00}\dot\varphi^2
   + g^{00}\dot\varphi\,\delta\dot\varphi + V'\delta\varphi - \frac13\,\delta\bar{{\cal L}}_r\right)\\
&\  + \frac13\, (A_r-1) \,\nabla_j\,\delta\bar{{\cal L}}_r
+\frac13\,\rho_{r}\, \del_{\varphi} \ln A_r \,\partial_j\delta\varphi \ .\nn
\end{align}
Once again, we distinguished the novelties due to the couplings.

\section{Fluid equations used for implementation}

The perturbation equations above imply the following source terms for the radiation sector
alone,
\beq
\nabla_{\mu} T^{\mu}{}_{\nu(r)}
=- \rho_r\,\del_\varphi \ln A_r\,\del_\nu\varphi
=-\bar\rho_r\,\del_\varphi A_r\,\del_\nu\varphi \ .
\eeq
This sign follows from the scalar equation of motion,
\beq
\ddot\varphi+3H\dot\varphi+V' +\bar\rho_m\,\del_\varphi A_m
+\bar\rho_r\,\del_\varphi A_r=0 \ ,
\eeq
and total energy-momentum conservation. At the background level,
\beq
\dot\rho_r+4H\rho_r=\rho_r\,\del_\varphi\ln A_r\,\dot\varphi \ ,
\eeq
which is also obtained directly by differentiating $\rho_r=A_r\bar\rho_r$ and
$\bar\rho_r\propto a^{-4}$.

For CLASS, using conformal time, $\mathcal{H}=a'/a$, and full coupled perturbation
variables
\beq
\delta_r\equiv\frac{\delta\rho_r}{\rho_r}\ ,\qquad
\theta_r=\bar\theta_r\ ,\qquad
\delta_r=\bar\delta_r+\Lambda_r\,\delta\varphi\ ,
\eeq
with
\beq
\Lambda_r\equiv \del_\varphi\ln A_r=\frac{A_r'}{A_r}\ ,\qquad
\Lambda_{2,r}\equiv \del_\varphi^2\ln A_r=\frac{A_r''}{A_r}-\Lambda_r^2\ ,
\eeq
the continuity and Euler equations are
\begin{align}
\delta_r' &=
-\frac43\left(\theta_r+\frac{h'}{2}\right)
+\Lambda_r\,\delta\varphi'
+\Lambda_{2,r}\,\varphi'\,\delta\varphi \ ,
\\
\theta_r' &=
\frac{k^2}{4}\delta_r-k^2\sigma_r
+\frac34\,\Lambda_r k^2\delta\varphi
-\Lambda_r\varphi'\theta_r \ .
\end{align}
Equivalently, in barred variables,
\begin{align}
\bar\delta_r' &=-\frac43\left(\bar\theta_r+\frac{h'}{2}\right) \ ,
\\
\bar\theta_r' &=\frac{k^2}{4}\bar\delta_r-k^2\bar\sigma_r
+\Lambda_r k^2\delta\varphi-\Lambda_r\varphi'\bar\theta_r \ .
\end{align}
The factor $3/4$ in the full-variable Euler equation is the force per inertial density,
since $\rho_r+p_r=\frac43\rho_r$; after changing to barred variables it combines with
$\delta_r=\bar\delta_r+\Lambda_r\delta\varphi$ to give the unit coefficient above.

The scalar perturbation equation receives the corresponding matter and radiation sources,
\beq
\delta\varphi''+2\mathcal{H}\delta\varphi'
+\left(k^2+a^2V''\right)\delta\varphi
+\frac12 h'\varphi'
=-a^2\sum_{s=m,r}\rho_s
\left(\Lambda_s\delta_s+\Lambda_{2,s}\delta\varphi\right) \ ,
\eeq
where the displayed equation uses physical densities. If CLASS variables store
$\rho_{\rm class}=\rho_{\rm phys}/3$, the right-hand side must be multiplied by the
corresponding factor of $3$. In barred variables, the radiation contribution is equivalently
$-a^2\left(A_r'\bar\rho_r\bar\delta_r+A_r''\bar\rho_r\delta\varphi\right)$.

The functions $A_m(\varphi)$ and $A_r(\varphi)$ should remain free implementation
parameters. The two numerical choices in Eqs.~\eqref{model1}--\eqref{model2} are benchmark
models for plots and comparisons.

%\newpage

$$$$

\addcontentsline{toc}{section}{References}
\bibliographystyle{JHEP}
\bibliography{References}


\end{document}  
